\documentclass[11pt,a4paper]{ctexart}
\usepackage[margin=2.6cm]{geometry}
\usepackage{amsmath,amssymb,bm}
\usepackage{booktabs}
\usepackage{enumitem}
\usepackage{setspace}
\usepackage{xcolor}
\usepackage{hyperref}
\usepackage{tcolorbox}
\tcbuselibrary{breakable}
\usepackage{fancyhdr}
\usepackage{listings}
\usepackage{microtype}

\hypersetup{colorlinks=true,linkcolor=blue!50!black,urlcolor=blue!60!black}
\definecolor{lockfr}{HTML}{2E7D32}
\definecolor{lockbg}{HTML}{E8F5E9}
\definecolor{warnfr}{HTML}{C62828}
\definecolor{warnbg}{HTML}{FFEBEE}
\definecolor{notefr}{HTML}{1565C0}
\definecolor{notebg}{HTML}{E3F2FD}
\newtcolorbox{notebox}{colback=notebg,colframe=notefr,title={说明},fonttitle=\bfseries,breakable,sharp corners=southwest}
\newtcolorbox{warnbox}{colback=warnbg,colframe=warnfr,title={训练合同},fonttitle=\bfseries,breakable,sharp corners=southwest}
\lstset{
  basicstyle=\ttfamily\small,
  breaklines=true,
  frame=single,
  columns=fullflexible,
  keepspaces=true,
  backgroundcolor=\color{gray!6}
}

\setlength{\headheight}{14pt}
\pagestyle{fancy}
\fancyhf{}
\lhead{TabUF Episode API 接口文档}
\rhead{v0 \quad 2026-08-14}
\cfoot{\thepage}
\setstretch{1.15}

\title{TabUF Episode API 接口文档\\[0.3em]\large 第零版观测数据服务}
\author{仓库 \texttt{1587causalai/tabuf-episode-api}}
\date{2026 年 8 月 14 日}

\begin{document}
\maketitle
\tableofcontents
\newpage

\section{这份文档管什么}

生成语义（总体、unit-specific 响应律、C/Q 出题）写在
\texttt{docs/data-generation.pdf}。
\textbf{这份文档只写运输层}：地址、方法、字段、状态码、例子。
字段名可以改，改的时候先改生成文档，再改这里，再改代码。

交互式文档由 FastAPI 自动生成，和本文同步：

\begin{itemize}
  \item Swagger UI：\texttt{/docs}
  \item ReDoc：\texttt{/redoc}
  \item OpenAPI JSON：\texttt{/openapi.json}
\end{itemize}

\section{服务地址}

当前部署在文献调研这台机器上，本机回环 + 一条临时公网隧道。

\begin{center}
\begin{tabular}{@{}ll@{}}
\toprule
角色 & URL \\
\midrule
本机 & \url{http://127.0.0.1:8787} \\
公网隧道（会变） & \url{https://spotlight-predicted-heaven-clips.trycloudflare.com} \\
\bottomrule
\end{tabular}
\end{center}

\begin{notebox}
隧道是 Cloudflare Quick Tunnel，没有 SLA，URL 在进程重启后会换。
稳定入口始终是本机 \texttt{127.0.0.1:8787}。
无鉴权、无 HTTPS 终止以外的证书、无速率配额（单次最多 32 episode，格子上限见下表）。
\end{notebox}

\section{约定}

\begin{itemize}
  \item 协议：HTTP/1.1，JSON，UTF-8。
  \item \texttt{Content-Type: application/json}。
  \item 无认证头。
  \item 时间：无时钟字段。可复现性靠 \texttt{seed}。
  \item 浮点：IEEE-754 JSON number。查询格在 \texttt{y\_input} 里是 JSON \texttt{null}，不是字符串 \texttt{"NaN"}。
  \item 掩码：JSON boolean。
  \item 拉直顺序：行优先（C-order）。\texttt{y\_query[t]} 对应 \texttt{query\_mask} 拉直后第 $t$ 个 \texttt{true}。
\end{itemize}

\section{端点一览}

\begin{center}
\begin{tabular}{@{}llll@{}}
\toprule
方法 & 路径 & 作用 & 鉴权 \\
\midrule
\texttt{GET}  & \texttt{/health} & 存活 & 无 \\
\texttt{GET}  & \texttt{/docs} & Swagger UI & 无 \\
\texttt{GET}  & \texttt{/redoc} & ReDoc & 无 \\
\texttt{GET}  & \texttt{/openapi.json} & 机器可读 schema & 无 \\
\texttt{POST} & \texttt{/v0/episodes} & 抽取观测 episode & 无 \\
\bottomrule
\end{tabular}
\end{center}

URL 里的 \texttt{/v0} 是生成语义版本，不是 HTTP 框架版本。
换响应律或换张量合同，才升 \texttt{/v1}。

\section{\texttt{GET /health}}

探活。生成器能 import、进程在听端口，就返回 200。

\paragraph{响应 200}
\begin{lstlisting}
{"ok": true, "version": "v0"}
\end{lstlisting}

\begin{center}
\begin{tabular}{@{}llp{7.5cm}@{}}
\toprule
字段 & 类型 & 含义 \\
\midrule
\texttt{ok} & bool & 恒为 true；能返回就说明活着 \\
\texttt{version} & string & 当前生成语义版本，现为 \texttt{v0} \\
\bottomrule
\end{tabular}
\end{center}

\section{\texttt{POST /v0/episodes}}

抽一批观测 episode。每集独立：先抽总体 $\{\mathbf{u}_i\}$，再填
$Y_{ij}=\langle\mathbf{u}_i,\mathbf{w}_j\rangle+e_{ij}$，再切 $C/Q$。

\subsection{请求体}

全部字段可选。缺省等于下表默认值。非法值返回 422。

\begin{center}
\begin{tabular}{@{}llclp{5.6cm}@{}}
\toprule
字段 & 类型 & 默认 & 范围 & 含义 \\
\midrule
\texttt{n\_units} & int & 64 & $[8,512]$ & 本集 unit 数 $n$ \\
\texttt{n\_features} & int & 8 & $[2,64]$ & 列数 $d$ \\
\texttt{unit\_dim} & int & 4 & $[1,32]$ & 个体潜维 $k$ \\
\texttt{query\_frac} & float & 0.15 & $(0,1)$ & $Q$ 占总格子比例 \\
\texttt{sigma} & float & 0.3 & $[0,10]$ & 观测噪声标准差 \\
\texttt{seed} & int/null & 0 & --- & 基种子；第 $e$ 集用 $\mathrm{seed}+e$ \\
\texttt{n\_episodes} & int & 1 & $[1,32]$ & 一次返回多少集（硬顶 32） \\
\texttt{debug} & bool & false & --- & 是否附带 $U,W,Y_{\mathrm{full}}$ \\
\bottomrule
\end{tabular}
\end{center}

\texttt{seed=null} 时服务端随机抽一个基种子，仍然写入每集的 \texttt{seed} 字段，所以返回值可复现、请求不必可复现。

最小请求：空对象 \texttt{\{\}}，等价于 $64\times 8$、$k=4$、$q=0.15$、1 集。

\subsection{响应 200}

\begin{lstlisting}
{
  "episodes": [
    {
      "context_mask": [[true, false], ...],
      "query_mask":   [[false, true], ...],
      "y_input":      [[0.12, null], ...],
      "y_query":      [0.87, ...],
      "shapes": {
        "n_units": 16, "n_features": 4, "unit_dim": 4,
        "n_query": 10, "n_context": 54,
        "y_input": [16, 4],
        "context_mask": [16, 4],
        "query_mask": [16, 4],
        "y_query": [10]
      },
      "seed": 0
    }
  ]
}
\end{lstlisting}

\begin{center}
\begin{tabular}{@{}llp{7.2cm}@{}}
\toprule
字段 & 形状 & 含义 \\
\midrule
\texttt{context\_mask} & $n\times d$ bool & $M^C$；true $\Leftrightarrow$ 该格 $\in C$ \\
\texttt{query\_mask} & $n\times d$ bool & $M^Q$；true $\Leftrightarrow$ 该格 $\in Q$ \\
\texttt{y\_input} & $n\times d$ float/null & $C$ 上为 $Y_{ij}$，$Q$ 上为 \texttt{null} \\
\texttt{y\_query} & $n_Q$ float & $Q$ 上真值，行优先拉直 \\
\texttt{shapes} & 对象 & 以上张量形状，避免客户端猜 \\
\texttt{seed} & int & 本集实际使用的种子 \\
\texttt{U} & $n\times k$ & 仅 \texttt{debug=true} \\
\texttt{W} & $d\times k$ & 仅 \texttt{debug=true} \\
\texttt{Y\_full} & $n\times d$ & 仅 \texttt{debug=true} \\
\bottomrule
\end{tabular}
\end{center}

不变量（第零版）：

\begin{align*}
  C\cap Q&=\emptyset,\\
  C\cup Q&=\{1,\ldots,n\}\times\{1,\ldots,d\},\\
  n_Q&=\mathrm{clip}(\mathrm{round}(q\cdot nd),1,nd-1).
\end{align*}

默认 payload \textbf{不含} $U,W,Y_{\mathrm{full}}$。
\texttt{debug=true} 时才出现这三个键；\texttt{false} 时键直接省略，不是 \texttt{null}。

\begin{warnbox}
训练 DataLoader 禁止设 \texttt{debug=true}，禁止把 $U$ 喂进网络。
$u_i$ 是个体真名。漏出去，预训练就是漏题。
\end{warnbox}

\subsection{错误}

\begin{center}
\begin{tabular}{@{}clp{8cm}@{}}
\toprule
状态码 & 何时 & 体 \\
\midrule
200 & 生成成功 & 见上 \\
422 & 字段越界、类型不对、缺 JSON & FastAPI 校验错误列表 \\
405 & 对 \texttt{/v0/episodes} 用了 GET 等 & --- \\
404 & 路径不存在 & --- \\
\bottomrule
\end{tabular}
\end{center}

第零版没有业务错误码。格子太大被 \texttt{le=} 拦在 422，而不是生成到一半再 500。

\section{调用例子}

探活：
\begin{lstlisting}
curl -s http://127.0.0.1:8787/health
\end{lstlisting}

最小生成：
\begin{lstlisting}
curl -s -X POST http://127.0.0.1:8787/v0/episodes \
  -H 'Content-Type: application/json' \
  -d '{}'
\end{lstlisting}

指定格子（推荐写进脚本的默认）：
\begin{lstlisting}
curl -s -X POST http://127.0.0.1:8787/v0/episodes \
  -H 'Content-Type: application/json' \
  -d '{
        "n_units": 16,
        "n_features": 4,
        "unit_dim": 4,
        "query_frac": 0.15,
        "sigma": 0.3,
        "seed": 0,
        "n_episodes": 2,
        "debug": false
      }'
\end{lstlisting}

公网隧道把主机换成当前隧道 URL 即可，路径不变。

\section{和生成文档的分工}

\begin{center}
\begin{tabular}{@{}ll@{}}
\toprule
问题 & 去哪看 \\
\midrule
为什么先抽总体 & \texttt{docs/data-generation.pdf} \\
$Y_{ij}=\langle u_i,w_j\rangle+e_{ij}$ & 同上 \\
$C/Q$ 为什么是格子级 & 同上 \\
这个字段叫什么、范围多少 & \textbf{本文} \\
现在服务在哪 & 本文第 2 节 \\
\bottomrule
\end{tabular}
\end{center}

接口改了，先改本文表格，再改 \texttt{app.py} 的 Pydantic 约束，最后才 bump 例子。

\end{document}
